\subsubsection{Datasets}
\label{subsubsec:datasets}

Er bestaan verschillende datasets die worden gebruikt om modellen te trainen en de prestaties van detectiesystemen te evalueren \autocite{Salloum2022}. Zo introduceerden \textcite{Mendes2025} MeAJOR, een verzameling van bestaande e-mailcorpora die onder meer bestaat uit de Enron-corpus \autocite{klimt2004}, met voornamelijk legitieme e-mails, en de Nazario-corpus \footfullcite{Nazario2025}, die phishingmails bevat.

Er zijn volgens \textcite{Pajola2025} wel enkele beperkingen aan de bestaande corpora:
\begin{enumerate}
    \item \textbf{Verouderde data.} Heel wat corpora kennen hun oorsprong van voor 2010 en capteren daarom de huidige aanvalstechnieken onvoldoende. 
    \item \textbf{Taal.} De meeste corpora bevatten engelstalige e-mails, wat de bruikbaarheid voor andere talen beperkt.
    \item \textbf{Ongelijkvormige data.} Niet elk corpus bewaart dezelfde (meta)data, waardoor kenmerken die headers of HTML vereisen niet in elke dataset kunnen worden geëxtraheerd.
\end{enumerate}
